\documentclass[12pt]{article}
\usepackage[utf8]{inputenc}
\usepackage[T1]{fontenc}
\usepackage{amsmath,amssymb}
\usepackage{geometry}
\geometry{margin=2.5cm}

\begin{document}

\begin{center}
\fbox{\textbf{Teorema (Cauchy)}}
\end{center}

\bigskip
\noindent Dacă $G$ este un gr.\ abelian finit şi $p\mid$ ord $G$, $p$=prim

\noindent Atunci $\exists\, a\in G$, ord $a=p$.

\bigskip
\noindent\underline{Dem}:

\noindent Not $n=$ ord $G$, $n\geq p$.

\noindent inducție după $n$.

\bigskip
\noindent $n=p=$ ord $G$.

\noindent $p$ prim $\Rightarrow p>1 \Rightarrow \exists\, a\neq e$. \ ord $a\mid$ ord $G=p$.

\noindent $a\neq e \Rightarrow$ ord $a>1$

\noindent $p=$prim

\noindent $\Rightarrow$ ord $a=p$.

\bigskip
\noindent\big(ord $G=p$=prim $\Rightarrow (\forall)\,a\in G\setminus\{e\}$ avem ord$(a)=p$\big).

\bigskip
\noindent pp $n>p$, $p\mid n$, şi că afirmația e adevărată pt.\ gr.\ de ordin $<n$.

\bigskip
\noindent\underline{Cazul I}: \ $\exists\, a\in G$ a.î.\ $p\mid$ ord$(a)$. \ $\Rightarrow$ ord$(a)=p\cdot l$.\footnote{simbolul exact folosit pentru acest exponent e neclar în original -- din raționamentul matematic trebuie să fie $l=\text{ord}(a)/p$.}

\noindent Iar $b=a^l \Rightarrow$ ord $b=p$.

\bigskip
\noindent\underline{Cazul II}: \ $\nexists\, a\in G$ a.î.\ $p\mid$ ord$(a)$.

\bigskip
\noindent Fie $a\in G$, $a\neq e$. \ pp.\ ord$(a)=m \Rightarrow p\nmid m$

\noindent Fie $H=\{e,\ldots,a^{m-1}\}\trianglelefteq G$ \ (pt.\ că $G$ e abelian).

\noindent Fie $\dfrac{G}{H}$. \ $\Rightarrow \left|\dfrac{G}{H}\right| = [G:H] < n$ \ (din Lagrange)

\bigskip
\noindent ord $G = n = m\cdot[G:H]$.

\[
p\mid n = m\cdot[G:H]
\]
\[
p\nmid m
\]
\[
\Rightarrow p\mid [G:H] = \left|\dfrac{G}{H}\right| < n.
\]

\bigskip
\noindent\underline{ip.\ de inducție} $\Rightarrow \exists\, \hat b\in \dfrac{G}{H}$ a.î.\ ord $\hat b=p$; \ $b\in G$=finit.\ $G$

\[
\Rightarrow \text{ord } b<\infty.
\]
\[
b^t=e
\]
\[
\hat e=\widehat{b^t}=\hat b^t \Rightarrow p\mid t.
\]

\noindent\textit{cu ip.\ de la cazul II.}\footnote{semnul de la finalul paginii pare mai degrabă un marcaj de tipul ,,q.e.d.'' decât o contradicție: $p\mid t=\text{ord}(b)$ e exact faptul care încheie Cazul II, reducându-l la tehnica deja stabilită în Cazul I (elementul $b^{t/p}$ ar avea ordin $p$).}

\end{document}
